\section{Introduction}
\label{sec:introduction}

The Zariski cancellation problem asks whether an affine variety is determined
by its cylinder when that cylinder is affine space. Over the complex numbers,
its algebraic formulation is the implication
\begin{equation}
 A[w]\simeq \mathbb C\poly{n+1}
 \quad\Longrightarrow\quad
 A\simeq\mathbb C\poly n,
 \label{eq:cancellation-question}
\end{equation}
where $A$ is a finitely generated commutative unital complex algebra, $w$ is
an independent polynomial variable, and all isomorphisms fix the complex
scalars. Here $C\poly r$ denotes a polynomial ring in $r$ variables over a
ring $C$. The hypothesis says that adjoining one variable turns $A$ into a
polynomial ring; the question is whether that variable can be cancelled.
An isomorphism of cylinders need not respect the given projection to the
affine line, so its existence alone does not identify polynomial
coordinates on $A$.

Cancellation for the affine line follows from the work of Abhyankar,
Heinzer, and Eakin~\cite[Corollary~(2.8)]{AEH}. Cancellation for the
complex affine plane was established through the work of Fujita,
Miyanishi, and Sugie~\cite{Fujita,MiyanishiSugie}, using the structure of
affine surfaces and cylinderlike open sets. In positive characteristic,
Asanuma's work on polynomial fiber rings supplied threefolds whose
one-variable cylinders are affine four-space~\cite{Asanuma}. Gupta
proved that suitable members of this family are not affine three-space
over their ground field~\cite{GuptaThreefold}, then established failure
of cancellation in every dimension at least three over every field of
positive characteristic~\cite{GuptaPositive}. The latter result uses a
separate family of examples: adjoining variables to a threefold whose
first cylinder is already polynomial would not produce higher-dimensional
counterexamples.

The characteristic-zero problem in dimensions at least three is still
listed as open in Gaifullin--Petrov's July 2026
preprint~\cite[p.~1]{GaifullinPetrov}. Gupta's
survey~\cite{GuptaSurvey} explains its history and its connections with
coordinates, polynomial fibrations, and the recognition of affine
space. We give a negative answer to \eqref{eq:cancellation-question}
over $\mathbb C$ for $n=4$.
The defining polynomial also gives a negative answer to the Stable
Coordinate Conjecture over $\kk$ in ambient dimension five, already after
adjoining one variable.

Let $p,s,u,F,J$ be algebraically independent over $\kk$, and set
\begin{equation}
 \begin{split}
 P&=\kk[p,s,u,F,J],\qquad x=s^2+u^3+p^2F,\\
 H&=x^2F-(1+2sx)J-p^2J^2-pu,\qquad A=P/(H).
 \end{split}
 \label{eq:example}
\end{equation}
Thus $x$ is an abbreviation for the displayed polynomial, and $H$ is a
polynomial in exactly the five indicated variables.

\begin{theorem}
\label{thm:main}
The algebra $A$ in \eqref{eq:example} is an integral complex algebra of
dimension four. For an independent variable $w$,
\[
 A[w]\simeq\kk\poly5,
 \qquad
 A\not\simeq\kk\poly4.
\]
Consequently the Zariski cancellation problem over the complex numbers has
a negative answer in dimension four.
\end{theorem}

\paragraph{The polynomial cylinder.}
The stabilization belongs to the theory of residual and stable
coordinates. A polynomial in $\kk[p]\poly4$ is a \emph{residual
coordinate} if its image is a polynomial coordinate over the residue
field of every prime of $\kk[p]$; it is \emph{one-stably a coordinate}
if it becomes a coordinate after adjoining one variable. Our $H$ is a
residual coordinate: it is a coordinate after inverting $p$, while a
linear change in $F,J$ makes its reduction at $p=0$ a coordinate.
Its one-stability therefore follows from the theorem of
El Kahoui--Ouali~\cite[Theorem~1.1]{ElKahouiOuali}, generalized by
Dutta--Lahiri~\cite[Theorem~3.4]{DuttaLahiri}.

The Stable Coordinate Conjecture asks whether a polynomial
$f\in\kk\poly n$ that becomes a coordinate in $\kk\poly{n+1}$ after
adjoining one variable must already be a coordinate in $\kk\poly n$
\cite[p.~342]{MakarLimanovVanRossumShpilrainYu}. The assertion concerns
polynomial automorphisms of the ambient spaces.

\begin{corollary}[Failure of the Stable Coordinate Conjecture]
\label{cor:stable-coordinate}
The polynomial $H$ in \eqref{eq:example} is not a coordinate of
$P=\kk[p,s,u,F,J]$, but it is a coordinate of $P[w]$ for one independent
variable $w$. Thus the Stable Coordinate Conjecture over $\kk$ fails in
ambient dimension five after one stabilization.
\end{corollary}

\begin{proof}
Since $H$ is a residual coordinate of $P=\kk[p]\poly4$, the theorem of
Dutta--Lahiri~\cite[Theorem~3.4]{DuttaLahiri}, applied with $R=\kk[p]$
and four relative variables, gives $G_2,G_3,G_4,G_5\in P[w]$ such that
\[
 P[w]=\kk[p][H,G_2,G_3,G_4,G_5]
     =\kk[p,H,G_2,G_3,G_4,G_5].
\]
Hence $H$ is a coordinate of $P[w]$ as a $\kk$-algebra. If $H$ were a
coordinate of $P$ as a $\kk$-algebra, quotienting a polynomial
coordinate system containing $H$ by $(H)$ would give
$A=P/(H)\simeq\kk\poly4$, contrary to \Cref{thm:main}.
\end{proof}

A derivation $D$ is \emph{locally nilpotent} if every element is killed
by some iterate of $D$; over $\kk$, its exponential is an algebraic
action of the additive group $\Ga$. We make the stabilization method
explicit. A locally nilpotent derivation of $P$
sends $H$ to $p^3$. Exponentiating it converts $H$ into $H+p^3w$,
and a coordinate calculation identifies the resulting quotient with a
polynomial ring. This exponential-shift argument follows the method of
El Kahoui--Ouali~\cite[Lemma~3.1 and proof of Theorem~1.1]{ElKahouiOuali};
the final coordinate change realizes the transfer criterion of
Edo--V\'en\'ereau~\cite[Theorem~7]{EdoVenereau}. We include its
polynomial inverse and the divisibility that makes it valid at $p=0$.

\paragraph{Affine-space fibrations.}
An $\Aff^3$-fibration is a morphism of complex affine varieties whose
base change to every geometric point of the target is affine three-space.
The affine-fibration assertion of Dolgachev and Weisfeiler predicts that
an affine morphism to a normal locally noetherian integral base is
Zariski-locally a product with affine $r$-space if every fiber is affine
$r$-space over its residue field~\cite[Section~3.8.5]{DolgachevWeisfeiler}.
The residual-coordinate property of $H$ gives two counterexamples in
relative dimension three:
\[
 p:\Spec A\longrightarrow\Aff^1,
 \qquad
 (p,H):\Spec P=\Aff^5\longrightarrow\Aff^2.
\]
Both maps are smooth $\Aff^3$-fibrations, but neither is a
Zariski-locally trivial $\Aff^3$-bundle. \Cref{sec:affine-fibrations}
proves this consequence using the locally-polynomial-algebra theorem
of Bass--Connell--Wright~\cite[Theorem~3(a)]{BassConnellWright} and the
Quillen--Suslin theorem~\cite[Theorem~2(b)]{Suslin}.

\paragraph{From polynomiality to a homogeneous additive action.}
The main work is to prove that $A$ itself is not polynomial.
Additive actions have long served to distinguish affine varieties from
affine space: Makar-Limanov used them to show that the Russell cubic is
not affine three-space~\cite{MakarLimanov}. Our obstruction has three
stages, beginning with the passage from filtered derivations to
homogeneous derivations developed by Kaliman and
Makar-Limanov~\cite[Sections~3--5]{KalimanMakarLimanov}.

The construction in \Cref{sec:construction,sec:degeneration} identifies
an auxiliary subalgebra of $A$ with the presentation
\[
 S=\kk[p,x,y,z,u]/\bigl(xy-z(z+1)-p^3u\bigr),
 \qquad S[p^{-1}]=A[p^{-1}].
\]
Under this inclusion, $p^2F$ and $p^2J$ belong to $S$. Define
\[
 R=S/pS=\kk[x,y,z,u]/\bigl(xy-z(z+1)\bigr),\qquad
 I=(\overline{p^2F},\overline{p^2J})\subset R,
\]
where the bars denote residues of elements of $S$ modulo $pS$.
The order $\ord_p$ along $p=0$ in $S$ applies to every nonzero
element of $A$. Its negative defines the \emph{valuation degree}
$\deg h=-\ord_p(h)$ and a filtration of $A$. We prove the exact
identification
\[
 G=\gr A=R[\tau,I/\tau^2]\subset R[\tau,\tau^{-1}],
 \qquad \deg\tau=-1,\qquad \deg R=0.
\]
Here $\tau$ is an indeterminate representing the initial form of $p$,
and the right-hand side is generated by $R$, $\tau$, and all
$i/\tau^2$ with $i\in I$. The coefficient inclusion is the one induced
by $S\subset A$: the residue of $b\in S$ corresponds to its
degree-zero class. This is an affine modification of $R[\tau]$ in
the sense of Kaliman--Zaidenberg~\cite{KalimanZaidenberg}, and it is
smooth.

If $A$ were polynomial, some polynomial coordinate would have positive
valuation degree: otherwise all of $A$ would have nonpositive degree,
contrary to $\deg F=2$. Differentiation in a different coordinate fixes
that positive-degree coordinate. The exact graded identification then
induces a nonzero homogeneous locally nilpotent derivation of $G$ with
a nonzero homogeneous invariant of positive degree. This is the action
the remaining two stages exclude.

\paragraph{Lifting through the principalizing line bundle.}
The standard torus bundle from $\mathrm{SL}_2$ to the smooth affine
quadric principalizes $I$ after adjoining the unchanged variable $u$.
This quadric and its polynomial embeddings also occur in the work of
Blanc and van Santen~\cite{BlancVanSanten}. The bundle algebra and
the pullback of $G$ are
\[
 \begin{gathered}
 \Rt=\kk[a,d,b,c,u]/(ac-bd-1),\qquad
 x=ab,\quad y=dc,\quad z=db,\\
 \Gt=G\otimes_R\Rt
     =\Rt[\tau,V]/(\tau^2V-v),\qquad
 v=a^3b-a^2u^3-d^2.
 \end{gathered}
\]
The valuation degree is zero on $\Rt$, with $\deg\tau=-1$ and
$\deg V=2$. Separately, fiber scaling gives $a,d$ weight $1$,
$b,c$ weight $-1$, and $u,\tau$ weight $0$; thus $v,V$ have
fiber weight $2$.

The morphism $\Spec\Gt\to\Spec G$ is the complement of the zero
section in a line bundle. Every additive action on the smooth affine
base $\Spec G$ lifts uniquely to an action on this line bundle that is
linear on fibers and
normalized at the identity. For the homogeneous action just obtained,
uniqueness makes the lift homogeneous with the same valuation degree
shift, while fiber linearity preserves fiber weights. This applies even
when the original action does not preserve the quadric algebra $R$.
A positive-degree invariant is divisible by $V$, forcing
$V$ itself to be invariant. The homogeneous pieces of $\Gt$, with
separate arguments according to whether $\tau$ is invariant, then
produce a nonzero locally nilpotent derivation $E$ of $\Rt$ with
nonpositive fiber weight and
\[
 E^2(v)=0.
\]
Thus the lift converts the geometric action into a constrained
derivation on the explicit determinant ring.

\paragraph{The weighted rigidity obstruction.}
A highest component $E'$ for an auxiliary grading remains a nonzero
locally nilpotent derivation with the same fiber weight. The grading
separates the two parts of $v$ so that $(E')^2(d^2+a^2u^3)=0$.
The polynomial orbit of $d^2+a^2u^3$ therefore has degree at most one.
The Mason--Stothers polynomial $abc$ inequality~\cite{Mason,Stothers}
and the nonsquareness of $-u^3$ in $\Frac\Rt$ force $E'$ to fix
$a,d,u$. The determinant relation then forces positive fiber weight,
contradicting the sign obtained in the preceding stage.
The use of polynomial $abc$ for rigidity of additive actions has
precedents in Finston--Maubach~\cite[Section~2]{FinstonMaubach}; here
the nonsquare obstruction requires keeping the leading coefficients
in the original fraction field.

\begin{table}[bp]
\centering
\small
\begin{tabular}{@{}llp{0.48\textwidth}@{}}
\toprule
Ring & Definition & Role in the proof\\
\midrule
$A$ & $P/(H)$ & Coordinate ring of the fourfold.\\
$S$ & \eqref{deg:S-R} & Subalgebra of $A$ defining the order along $p=0$.\\
$R$ & $S/pS$ & Coefficient ring of the degeneration.\\
$G$ & $\gr A=R[\tau,I/\tau^2]$ & Exact graded algebra carrying the induced action.\\
$\Rt$ & \eqref{bun:determinant-ring} & Torus-bundle algebra in which $I$ becomes principal.\\
$\Gt$ & $G\otimes_R\Rt$ & Pullback algebra carrying the lifted action.\\
\bottomrule
\end{tabular}
\caption{The rings in the proof. The valuation degree on $G$ and
$\Gt$ is separate from the fiber weight on $\Rt$ and $\Gt$.}
\label{tab:ring-dictionary}
\end{table}

\paragraph{Organization.}
\Cref{tab:ring-dictionary} records the rings used in the argument.
\Cref{sec:construction} proves the stabilization.
\Cref{sec:degeneration} constructs the filtration and identifies $G$.
\Cref{sec:bundle} describes the torus bundle and proves the lifting
statement. \Cref{sec:reduction} extracts the constrained derivation on
$\Rt$, and \Cref{sec:rigidity} excludes it and completes the proof of
\Cref{thm:main}. \Cref{sec:affine-fibrations} derives the two nontrivial
affine-space fibrations. All locally nilpotent derivations and algebraic group
actions in the proof are defined over $\kk$.
